\section{\texorpdfstring{Lipschitz 常数就是认证半径}{Lipschitz 常数就是认证半径}}
\label{sec:17-Constrained-Guarantees/02-Adversarial-Robustness/02}

一次安全评审上，问题只有一句：这个模型在 \(\norm{\delta}_2\le0.5\) 的扰动下会不会改变判断？

团队的回答是：我们跑了单步攻击、多步攻击、几种迁移攻击和一套公开的攻击集合，准确率从 \(96\%\) 掉到 \(71\%\)，没掉到零，所以模型是鲁棒的。

这个回答的问题不在数字，在逻辑。上一节末尾已经点明，鲁棒性是一个关于邻域的全称命题；而``跑了若干种攻击都没打穿''否定不了任何全称命题，它只说明这几种攻击没找到反例。明年出现一种新攻击，同一个模型可能就掉到零——这在过去十年反复发生过，一批当时看起来牢固的防御在新攻击面前逐个失守。

要真正回答那个问题，需要一个不依赖攻击方法的论证：给出一个半径 \(R\)，证明半径 \(R\) 的球内不存在任何反例。这个论证存在，而且短得出人意料。

\subsection{一个量控制着全部}

设分类器输出 \(K\) 个 logit \(f_1(x),\dots,f_K(x)\)，当前预测类别为 \(y\)。定义\textbf{间隔函数}
\[
g(x)=f_y(x)-\max_{y'\neq y}f_{y'}(x).
\]
\(g(x)>0\) 等价于预测为 \(y\)。要让预测在扰动下不变，只需保证 \(g\) 在整个邻域内不变号。

\begin{theorem}[Lipschitz 认证]
若 \(g\) 关于 \(\ell_2\) 范数是 \(L\)-Lipschitz 的，即对所有 \(u,v\) 有 \(\abs{g(u)-g(v)}\le L\norm{u-v}_2\)，则对一切 \(\norm{\delta}_2<g(x)/L\)，有 \(g(x+\delta)>0\)，即预测保持不变。
\end{theorem}

证明一行：
\[
g(x+\delta)\;\ge\;g(x)-\abs{g(x+\delta)-g(x)}\;\ge\;g(x)-L\norm{\delta}_2\;>\;g(x)-L\cdot\frac{g(x)}{L}=0 .
\]

这条结论与攻击方法无关，与优化器无关，与谁来验证无关。它是一个纯粹的分析陈述：函数变化的速度有上界，起点离零有距离，于是在一段距离内到不了零。半径 \(g(x)/L\) 叫做\textbf{认证半径}。

条件对账逐条来。\(L\) 必须是 \(g\) 的 Lipschitz 常数，而实践中通常估的是 \(f\) 作为向量值映射的常数；两者之间隔着一个 \(\sqrt2\) 量级的因子（取两个分量之差会放大），漏掉它会让声称的半径偏大。范数必须前后一致：用 \(\ell_2\) 的 Lipschitz 常数只能认证 \(\ell_2\) 球，要认证 \(\ell_\infty\) 球得换算，而换算会带来 \(\sqrt d\) 的损失——上一节讲过高维方块与球的比值就是 \(\sqrt d\)，这个因子在这里再次出现，而且是以最不利的方式。\(g(x)>0\) 是前提，模型本来就分错的点没有半径可言。最要紧的一条：\(L\) 必须是真正的\textbf{上界}。用有限个样本点上的差商去``估计''一个 Lipschitz 常数，得到的是下界，把下界代进定理就得到一个没有依据的半径，而这个错误不会以任何方式暴露出来。

\begin{center}
\begin{tikzpicture}[x=2.8cm,y=1.55cm,>=stealth,font=\small]
  \draw[->,black!60] (-0.12,0) -- (3.2,0) node[right,font=\footnotesize] {位移 \(t\)};
  \draw[->,black!60] (0,-0.9) -- (0,1.6) node[above,font=\footnotesize] {间隔 \(g\)};
  \draw[thick,blue!65] plot[domain=0:3.05,samples=90]
    (\x,{1.2-0.175*\x*\x+0.07*sin(\x*200)});
  \draw[very thick,orange!85!black] (0,1.2) -- (0.4,0);
  \draw[orange!85!black,dashed] (0.4,0) -- (0.53,-0.4);
  \draw[black!45,dashed] (0.4,0) -- (0.4,-0.3);
  \draw[black!45,dashed] (2.62,0) -- (2.62,-0.6);
  \fill[black] (0,1.2) circle (1.8pt);
  \node[anchor=east,font=\footnotesize] at (-0.03,1.2) {\(g(x)\)};
  \draw[<->,black!70] (0,-0.3) -- (0.4,-0.3);
  \node[anchor=west,font=\footnotesize,black!70] at (0.62,-0.3)
    {认证半径 \(g(x)/L\)};
  \draw[<->,black!70] (0,-0.6) -- (2.62,-0.6);
  \node[below,font=\footnotesize,black!70] at (1.35,-0.63) {真实的翻转位置};
  \node[anchor=west,font=\footnotesize,orange!85!black] at (0.5,0.34)
    {斜率 \(-L\) 的包络};
  \node[anchor=west,font=\footnotesize,blue!65] at (1.7,0.95) {真实的 \(g\)};
\end{tikzpicture}
\end{center}

\emph{要记住的是两个箭头的长度差：认证半径由``最坏可能的下降速度''推出，真实的翻转位置由函数实际的下降速度决定。包络越松，短箭头越短，而长箭头一动不动——保守的界不会犯错，只会浪费。}

\subsection{坏消息：这个常数算不出来}

定理把问题归结成了求 \(L\)。对线性模型，\(L=\norm{w}_2\)，一步就得。对神经网络，事情立刻变难。

精确计算一个 ReLU 网络的 Lipschitz 常数是 NP 难的，连两层网络都是。这不是算法还没找到，是这个问题落在\hyperref[sec:09-Convex-Optimization/09-Complexity-and-Approximation/01]{``P、NP 与 NP-hard：给问题按计算难度分类''}划出的那一类里，除非复杂度理论的基本猜想被推翻，否则不会有多项式算法。

实践中退而求其次，用一个可算的上界。复合函数的 Lipschitz 常数不超过各层常数之积，而线性层的 \(\ell_2\) Lipschitz 常数正是其权重矩阵的谱范数——最大奇异值，见\hyperref[sec:03-Linear-Algebra/07-SVD-and-Data-Applications/01]{``任意矩阵的奇异值分解''}。ReLU 与常见激活的常数是 1。于是
\[
L\;\le\;\prod_{\ell=1}^{m}\sigma_{\max}(W_\ell).
\]

这个界可以用幂迭代快速算出，问题是它松得厉害。松的原因有两处，都很具体。

一处是每层的 \(\sigma_{\max}\) 对应一个特定的输入方向，而相邻两层的最大奇异方向几乎不会对齐。第一层把某个方向放大 3 倍，第二层的 3 倍方向是另一个方向，实际的复合放大远小于 9。层数越多，这种方向失配累积得越厉害，乘积形式的上界因此随深度指数地偏离真实值。

另一处是激活函数的收缩被完全忽略。ReLU 在负半轴上把输入压成零，实际的局部斜率是 0 而不是 1；用 1 作为上界在数学上正确，在数量上极其浪费，而一个典型网络里相当比例的单元在任意给定输入上处于非激活状态。

两处相加，得到的上界比真实值大几个数量级是常态。带进定理，认证半径就被压缩到同样的倍数——一个实际能扛住 \(0.5\) 扰动的模型，谱范数乘积给出的认证半径可能只有 \(0.005\)。有更紧的界（半正定规划松弛、区间传播、以及两者的混合，思路上都属于\hyperref[sec:09-Convex-Optimization/09-Complexity-and-Approximation/02]{``凸松弛把 NP-hard 问题变成可解问题''}那一类），它们把差距缩小到一个数量级左右，换来的是每个输入点上要解一个优化问题，无法用于大规模评测。

\subsection{保守的界和乐观的估计，谁也替代不了谁}

现在手上有两个数：认证半径（下界，来自松的 Lipschitz 估计）和攻击成功的最小扰动（上界，来自若干次攻击尝试）。真实的鲁棒半径夹在两者之间，而这个区间往往跨两个数量级。

这两个数的\textbf{性质完全不同}，必须分开使用。

认证半径是保守的，但它是真的。说``在半径 \(0.005\) 内预测不变''，这句话对所有可能的扰动成立，包括还没被发明的攻击。它可能远小于实际能扛住的量，但它不会被推翻。

攻击给出的半径是乐观的，而且不可靠。说``现有攻击在 \(0.5\) 以内都没成功''，这句话只对试过的那些攻击成立。它更接近真实值，但它随时可能因为一种新攻击而作废，而且没有任何办法预先知道会不会作废。

工程上常见的错误是拿后者当结论、拿前者当``理论上过于保守可以忽略''。正确的读法是：前者是能承诺的，后者是能观察到的，中间那段区间是我们对这个模型的无知。想缩小这段区间只有两条路，把界做紧或者把攻击做强，两条都在进展，都没有终结。

这个区分与\hyperref[sec:11-Mathematical-Statistics/04-Interval-Estimation/01]{``置信度属于程序，而非这次参数''}里那条判断同源：一个保证的效力来自产生它的论证结构，不来自它给出的数字有多好看。攻击失败这件事没有对应的论证结构，所以它无论重复多少次都不升级成保证。

\subsection{随机平滑：绕开 Lipschitz 常数}

有一条路完全不去估 Lipschitz 常数，而是自己造一个 Lipschitz 常数已知的分类器。

给定任意基分类器 \(f\)（可以是任何东西，不要求可微、不要求已知结构），定义平滑分类器
\[
h(x)=\argmax_{c}\;P_{\eta\sim\mathcal N(0,\sigma^2 I)}\big(f(x+\eta)=c\big).
\]
即在输入上加各向同性高斯噪声，看基分类器最常输出哪一类。

关键结论是：设在 \(x\) 处最高的那一类概率为 \(p_A\)、次高为 \(p_B\)，则对一切
\[
\norm{\delta}_2<\frac{\sigma}{2}\Big(\Phi^{-1}(p_A)-\Phi^{-1}(p_B)\Big)
\]
有 \(h(x+\delta)=h(x)\)，其中 \(\Phi^{-1}\) 是标准正态的分位函数（见\hyperref[sec:05-Probability/03-Random-Variables/03]{``连续分布、密度与分位数''}）。证明的思路是：平移一个高斯分布，任何事件的概率变化都被平移量控制住，这个控制关系可以精确写出来。换句话说，高斯卷积把一个任意糟糕的函数变成了一个光滑程度可控的函数——Lipschitz 性质是被卷积\textbf{制造}出来的，不是被估计出来的。

这条路的好处很实在：它与网络的结构、深度、是否可微全都无关，因此可以用在大模型上；而且半径直接由 \(p_A\) 与 \(\sigma\) 给出，不经过任何松弛。

要付出的东西有三样，都不轻。

第一样是推理开销。\(p_A\) 无法解析计算，只能采样估计，通常要跑上千次前向（见\hyperref[sec:05-Probability/08-Limit-Theorems/05]{``Monte Carlo：从随机样本到数值答案''}）。一次预测的算力变成上千倍，这是实打实的计算开销。

第二样是保证的类型变了。用采样估的 \(p_A\) 只能给出一个置信下界，于是认证结论的形式是``以至少 \(1-\alpha\) 的概率，半径 \(R\) 内预测不变''。这不再是确定性的陈述，而是一个带显著性水平的统计陈述，形状与\hyperref[sec:11-Mathematical-Statistics/07-Statistical-Decision-and-Reliable-Prediction/05]{``Conformal prediction 用可交换性换有限样本覆盖''}给出的覆盖保证类似——概率属于认证程序，不属于这一次的判断。多个样本上反复调用同一套程序时，还要注意\hyperref[sec:11-Mathematical-Statistics/05-Hypothesis-Testing/04]{``多重比较改变错误的计量单位''}里的那条提醒。

第三样是 \(\sigma\) 的两难。半径与 \(\sigma\) 成正比，所以想认证大半径就得加大噪声；而噪声大了之后，基分类器在带噪输入上分不准，\(p_A\) 随之下降，分位函数那一项跟着缩水，半径未必真的变大。更直接的问题是平滑分类器本身的干净准确率随 \(\sigma\) 单调下降。每个 \(\sigma\) 对应一条``半径—准确率''曲线上的一个点，没有哪个 \(\sigma\) 在所有半径上都最优，实践中要对不同 \(\sigma\) 分别训练分别认证。

\begin{remark}
基分类器需要在带噪输入上训练才能有像样的 \(p_A\)，这属于经验做法而不是定理要求——认证结论对任何基分类器都成立，只是没经过噪声训练的基分类器给出的 \(p_A\) 太低，半径接近零。同样属于经验的还有采样数与 \(\alpha\) 的选取，公开工作里的取值范围很宽，没有理论指明最优点。
\end{remark}

\subsection{半径有了，接下来是准确率}

到这里，``鲁棒不鲁棒''这个问题有了可以负责的回答形式：给出一个半径，说明它是怎么算出来的，以及它是确定性的还是带置信水平的。

但这个回答还没触及一个更基本的问题。前面所有讨论都默认``鲁棒''和``准确''是可以同时改进的两件事，只是我们暂时不够聪明。下一节说明这个默认不成立：存在具体的数据分布，使得标准准确率与鲁棒准确率之和被一个小于 2 的常数卡住，于是一个上去另一个必然下来。那不是工程问题，是有构造性反例的结论。
